\scanpage{65}\mspage{64}

Notons que \uline{At} 9 est impliqué par chacun des deux \deleted{\unsure{A|a}}axiomes At 9$'$, At 9$''$. \deleted{Et il le \unsure{sera|serait} aussi \unsure{par} l'axiome de recollement plus bas, \unsure{donc il} n'est pas tellement crucial en lui-même.}
% slanted
\margin{\textcircled{$*$} \unsure{Intervertir|Intervention} At 9, At 10, pour \unsure{bloquer} les \unsure{ens} At 8, At 10 (qui sont \unsure{les|des} \unsure{2|7} axiomes \deleted{d'existence} portant sur les strates ouvertes : l'existence des lieux et des \unsure{0-sphères|o-sphères} \unsure{reliées|reliés} aux \unsure{strates} \unsure{ouvertes|issues de}\unsure{.|)}}
% slanted
\margin{\unsure{Donner|Donnons} ici deuxième Axiome de \unsure{disjonction} \uline{At 10} Soient $\deleted{\unsure{x|X}}X,\deleted{\unsure{y|Y}}Y\in M$ \inserted{$X \mathrel{\tikz[baseline=-.45ex,x=1em,y=1em]{\draw[line width=.055em] (.52,.27)--(.05,.055)--(.48,-.11); \draw[line width=.055em] (.14,.21)--(.46,.04)--(0,-.17);}} Y$} \deleted{une figure, $\unsure{x,y}\in\operatorname{omb}(F)$} $X\ne Y$, \inserted{alors} \deleted{\unsure{appartenant à|appart. à|portent} des \uline{strates} \uline{ouvertes} \uline{disjointes}} tout $x\in\operatorname{omb}(X)^{\unsure{\circ|0}}$ \deleted{\unsure{i.e.} $\unsure{\exists x_i|X_i}$} est \uline{disjoint} de tout $y\in\operatorname{omb}(Y)^{\circ}$.}
% slanted
\margin{\uline{NB} Le cas où \deleted{$Y\unsure{\leq|\nleq}X$ ou $X\leq Y$ est} $X\wedge Y\ne X,Y$ résulterait aussi de l'axiome de recollement At 11 plus bas -- donc le cas crucial est celui où $Y\subset X$.}
% slanted
\margin{Les axiomes At 8, At 9, At 10 s'appelleront ``axiomes des lieux''.}
% slanted
\margin{\uline{At 10} $X,Y\in M$, $X\ne Y$ \inserted{$x,y\in\unsure{L|\mathcal L|\alpha}$} $x\overset{\circ}{\ll}X,\ y\overset{\circ}{\ll}Y\ \Longrightarrow\ x\parallel y$ (ou encore $x \mathrel{\tikz[baseline=-.45ex,x=1em,y=1em]{\draw[line width=.055em] (.52,.27)--(.05,.055)--(.48,-.11); \draw[line width=.055em] (.14,.21)--(.46,.04)--(0,-.17);}} y$, cela revient au $=$) Voir (\unsure{ci corr.|ci-contre}) p 72 \textcircled{$*$}}

\uline{21. juin}

\uline{Proposition} Soient $\unsure{\mathcal F|F}$ un atelier At 1 -- \deleted{\unsure{At 8|At 9}} At 7 $F\in\mathcal F$, $F'\ll F$, $G\leq F$. Alors
\[
G'=\operatorname{Inf}^{\ll}(F',G)
\]
existe, et $G'$ est une sous-figure de $F'$ : \deleted{On}
% slanted
\margin{Peut remonter avant les histoires de lieux}
\diagram{\begin{tikzcd}[ampersand replacement=\&,row sep=2.25em,column sep=2.2em]
F' \arrow[r,phantom,"\ll" description] \arrow[d,phantom,"{\rotatebox{90}{$\leq$}}" description] \& F \arrow[d,phantom,"{\rotatebox{90}{$\leq$}}" description]\\
G' \arrow[r,phantom,"\ll" description] \& G.
\end{tikzcd}}

On a
\begingroup\makeatletter\tagsleft@true\makeatother
\[
\widetilde G'=\{X\in\widetilde F'\mid X\ll G\}\ \deleted{\unsure{=\operatorname{Inf}^{\leq}|=\operatorname{Inf}^{\ll}|\ll\operatorname{Inf}^{\leq}}}.\tag*{(1.60)}
\]
\endgroup
% slanted
\margin{Terminologie : \uline{Sous figure induite} $G'$ de $G$ par la sous-figure $F'$ de $F$. \uline{Raffinement} $G'$ de $G$, \uline{induit} sur la sous-figure $G$ de $F$ par le \unsure{raffinem.|raffinement} $F'$ de $F$, ou $F'/G$ \unsure{--|.}}

\uline{Démonstration} La formule (1.60) \deleted{résulte} \inserted{est} tautologique par la définition des $\operatorname{Inf}^{\ll}$ et du fait que $\leq$ implique $\ll$. Mais sans admettre l'existence du $\operatorname{Inf}$, la formule (1.60) définit une partie fermée de $\widetilde F'$, donc une sous-figure $G'$ de $F'$, et on a évidemment $G\ll F$ par
